\chapter{Точный двухузловой компактон}
\label{ch:v6-spectral-transition-discrete-compacton-existence-gate}

\section{Одноузловой запрет}

Для условного сдвига правая компонента любого узла переносится на соседний
узел справа, а левая --- на соседний узел слева. Локальная монета не может
создать амплитуду на других узлах до сдвига. Поэтому ненулевое состояние с
опорой на одном узле после шага имеет опору на двух соседях и не может
вернуться к себе с общей фазой.

Этот запрет не зависит от значения $\kappa$ и внутреннего состава:
одноузлового стационарного compacton нет.

\section{Двухузловая схема без утечки}

Пусть опора состоит из соседних узлов $0,1$. Чтобы после сдвига не было
утечки, на узле $0$ начальная амплитуда должна лежать только в левом
направлении и полностью переворачиваться в правое; на узле $1$ --- только
в правом направлении и полностью переворачиваться в левое.

Возьмём одинаковый сбалансированный хиральный вектор
\begin{equation}
 v=(\ell,e),\qquad
 \|\ell\|^2=|e|^2=\frac14,qquad \|v\|^2=\frac12.
 \label{eq:v6-compacton-local-vector}
\end{equation}
Для него собственная частота составного хирального генератора на активной
плоскости равна
\begin{equation}
 a=\|\ell\|\,|e|=\frac14.
\end{equation}
Полный разворот направления требует
\begin{equation}
 \cos(\kappa a)=0,
 \qquad
 \kappa a=\frac{(2m+1)\pi}{2}.
\end{equation}
Следовательно,
\begin{equation}
 \boxed{\kappa_m=2(2m+1)\pi},\qquad m=0,1,2,\ldots
 \label{eq:v6-compacton-kappa-series}
\end{equation}
и минимальная положительная ветвь равна $\kappa_0=2\pi$.

\section{Точное собственное состояние}

Обозначим два направления через $+$ и $-$. Для
$s_m=(-1)^m$ зададим
\begin{equation}
 \Psi_0=|-\rangle\otimes v,
 \qquad
 \Psi_1=|+\rangle\otimes i s_m v,
 \qquad
 \Psi_n=0\quad(n\ne0,1).
 \label{eq:v6-two-site-compacton}
\end{equation}
Локальная монета полностью меняет направления, после чего сдвиг меняет
узлы местами. Прямое вычисление даёт
\begin{equation}
 F_{\kappa_m}(\Psi)=i\Psi.
 \label{eq:v6-compacton-eigenphase}
\end{equation}
Комплексно сопряжённая ветвь имеет фазу $-i$. Таким образом, профиль имеет
точно два узла, период четыре без учёта общей фазы и норму один.

Машинный сертификат проверяет первые три ветви
$\kappa/\pi=2,6,10$. Максимальные остаток собственного состояния и утечка
меньше $10^{-15}$. После ста шагов минимальная ветвь сохраняет всю норму на
двух исходных узлах.

\section{Что именно квантовано}

Получено не физическое значение массы, а точное условие существования на
безразмерное произведение нелинейной связи и локальной нормы. При полной
нормировке симметричного двухузлового состояния оно превращается в
дискретный ряд \eqref{eq:v6-compacton-kappa-series}. Минимальность выделяет
$2\pi$, но пока отсутствуют:
\begin{enumerate}
 \item доказательство нелинейной устойчивости compacton;
 \item вывод физического шага решётки и размера $2a$;
 \item карта от $2\pi$ к слепой наблюдаемой Стандартной модели;
 \item механизм динамического захвата общего состояния в эту орбиту.
\end{enumerate}

Расстройка связи создаёт утечку уже за один шаг. Для
$\kappa=2\pi(1+\varepsilon)$ точная вероятность утечки равна
\begin{equation}
 P_{\rm leak}=\sin^2\!\left(\frac{\pi\varepsilon}{2}\right).
 \label{eq:v6-compacton-detuning-leakage}
\end{equation}
Это показывает, что существование дискретно, но ещё не отвечает на вопрос
об устойчивости относительно возмущений самого состояния.

\begin{theorem}[Существование минимального компакттона]
Минимальная составная пространственная монета не допускает ненулевого
одноузлового стационарного профиля, но допускает точный симметричный
двухузловой compacton. Его условие существования равно
$\kappa=2(2m+1)\pi$, собственная фаза равна $\pm i$, а минимальная
положительная ветвь даёт новое безразмерное значение $2\pi$.
\end{theorem}

\begin{nogocriterion}
Нельзя объявлять $2\pi$ предсказанной константой природы до доказательства
устойчивости и построения физической карты. Точная периодическая орбита
является кандидатом endpoint, но не доказанным аттрактором и не массовым
состоянием в физических единицах.
\end{nogocriterion}

\section{Следующий гейт}

Следующий гейт
\path{version6_spectral_transition_discrete_compacton_stability_quantization_gate}
должен вычислить Floquet/Jacobi-линеаризацию четырёхшаговой орбиты,
отделить калибровочные и фазовые нулевые направления и проверить рост
пространственных возмущений.

Нулевая гипотеза: компактон имеет неустойчивое направление утечки.
Стоп-критерий: спектральный радиус физической редуцированной монодромии
выше единицы закрывает его как устойчивую частицу; радиус не выше единицы
открывает вывод физического масштаба и карты значения $2\pi$.

\begin{protocol}
\begin{enumerate}
 \item одноузловой стационарный compacton: \textbf{запрещён};
 \item двухузловой compacton: \textbf{построен точно};
 \item условие существования: \textbf{$\kappa=2(2m+1)\pi$};
 \item минимальная ветвь: \textbf{$2\pi$};
 \item утечка точного решения: \textbf{ноль};
 \item фазовый период: \textbf{четыре шага};
 \item нелинейная устойчивость: \textbf{не доказана};
 \item физический размер и наблюдаемая карта: \textbf{не выведены};
 \item R4 и R5: \textbf{частично открыты}.
\end{enumerate}
\end{protocol}

Машинный сертификат сохранён в
\path{s2t_v6_spectral_transition_discrete_compacton_existence_gate_results.json}.